\documentclass{article}

\usepackage[utf8]{inputenc}
\usepackage[T2A]{fontenc}
\usepackage{helvet}
\usepackage{geometry}
\usepackage{xcolor}
\usepackage{eso-pic}
\usepackage{fancyhdr}
\usepackage{parskip}
\usepackage{microtype}
\usepackage[english,ukrainian]{babel}
\usepackage{url}
\usepackage{amsmath}
\usepackage{amssymb}
\usepackage{amsthm}
\usepackage{longtable}
\usepackage{booktabs}
\usepackage{array}
\usepackage{hyperref}
\usepackage{titlesec}
\usepackage{listings}
\usepackage{enumitem}

% Define brand colors
\definecolor{primary}{HTML}{293757}
\definecolor{footergray}{gray}{0.65}

\renewcommand{\familydefault}{\sfdefault}
\setcounter{secnumdepth}{3}

% Font shapes for T2A/Helvetica
\DeclareFontFamily{T2A}{phv}{}
\DeclareFontShape{T2A}{phv}{m}{n}{<->ssub * cmss/m/n}{}
\DeclareFontShape{T2A}{phv}{bx}{n}{<->ssub * cmss/bx/n}{}
\DeclareFontShape{T2A}{phv}{m}{it}{<->ssub * cmss/m/it}{}
\DeclareFontShape{T2A}{phv}{bx}{it}{<->ssub * cmss/bx/it}{}

\titleformat{\section} {\color{primary}\Huge\bfseries}{\thesection.}{1em}{} \titlespacing*{\section}{0pt}{30pt}{12pt}
\titleformat{\subsection} {\color{primary}\LARGE\bfseries}{\thesubsection.}{1em}{} \titlespacing*{\subsection}{0pt}{20pt}{8pt}
\titleformat{\subsubsection} {\color{primary}\Large\bfseries}{\thesubsubsection.}{1em}{} \titlespacing*{\subsubsection}{0pt}{10pt}{6pt}

\fancyhf{}
\fancyhead[R]{\small\color{footergray} 04 липня 2026}
\fancyfoot[L]{\small\color{footergray} Технічні рекомендації ERP/1. SRE Handbook}
\fancyfoot[R]{\small\color{footergray}\thepage}

\newcommand\HUGE[1]{\fontsize{38}{38}\selectfont #1}
\renewcommand{\headrulewidth}{0pt}
\renewcommand{\footrulewidth}{0.4pt}
\renewcommand{\footrule}{\color{footergray}\hrule width \textwidth height 0.4pt}

\lstset{
  basicstyle=\ttfamily\small,
  breaklines=true,
  frame=single,
  numbers=left
}

\begin{document}
\pagestyle{fancy}

\begin{titlepage}
\AddToShipoutPictureBG*{\AtPageLowerLeft{\color{primary}\rule{\paperwidth}{\paperheight}}}
\color{white}\thispagestyle{empty}\vspace*{4cm}\centering
  {\Huge \textbf{ERP/1: Підприємство} \par} \vspace{2cm}
  {\HUGE \textbf{SRE Handbook} \par} \vspace{2.5cm}
  {\LARGE \textbf{Практики Site Reliability Engineering \\ для Continuous Delivery ERP/1} \par} \vspace{2.5cm}
\vfill
  {\large 2026 \copyright\ Системи Електронного Урядування \par}
\end{titlepage}

\newpage
\begin{center}
  {\Huge\color{primary}\bfseries Зміст\par}
\end{center}
\vspace{40pt}
\tableofcontents

\newpage
\begin{abstract}
У цьому документі наведено практичний посібник з Site Reliability Engineering (SRE) для ERP/1: Continuous Development. Він адаптує класичні принципи SRE від Google до контексту Kubernetes-based, Helm-managed, multi-namespace розгортань ERP.uno. Посібник охоплює SLO/SLI, error budgets, on-call, monitoring, automation, chaos engineering та інтеграцію з Helm/GitOps для забезпечення високої надійності державних та корпоративних систем.
\end{abstract}

\newpage
\section{Умовні скорочення та визначення}

\begin{longtable}{p{4cm}p{10cm}}
\toprule
\textbf{Термін / Скорочення} & \textbf{Значення} \\
\midrule
SRE & Site Reliability Engineering (Інженерія надійності сайтів) \\
SLO & Service Level Objective (Цільовий рівень обслуговування) \\
SLI & Service Level Indicator (Індикатор рівня обслуговування) \\
SLA & Service Level Agreement (Угода про рівень обслуговування) \\
TOIL & Time spent On Internal Labor (Ручна робота) \\
K8s & Kubernetes \\
Helm & Package manager for Kubernetes \\
GitOps & Declarative continuous delivery \\
ERP/1 & Електронна система планування ресурсів версії 1 \\
CD & Continuous Development / Continuous Delivery \\
 observability & Спостережуваність (metrics, logs, traces) \\
Chaos Engineering & Інженерія хаосу (тестування відмовостійкості) \\
\bottomrule
\end{longtable}

\newpage
\section{Загальні відомості}

\subsection{Передумови}

Сучасні державні та корпоративні системи ERP/1 вимагають високої доступності, швидкого відновлення та автоматизованої надійності. Репозиторій \texttt{erpuno/cd} надає Helm-чарти, shared resources та lib для multi-namespace deployment'ів, що створює основу для SRE-практик.

Ключові виклики:
\begin{enumerate}
\item Забезпечення 99.9\%+ uptime в умовах потенційних кіберзагроз та інфраструктурних збоїв.
\item Мінімізація TOIL через автоматизацію GitOps.
\item Комплаєнс з КСЗІ та міжнародними стандартами надійності.
\item Масштабування observability для сотень вузлів.
\end{enumerate}

\subsection{Призначення SRE Handbook}

Цей handbook визначає:
\begin{itemize}
    \item Принципи SRE, адаптовані до ERP/1 CD pipeline.
    \item Конкретні SLI/SLO для namespaces (\texttt{erp-ai}, \texttt{erp-telemetry}, \texttt{erp-apps} тощо).
    \item Практики on-call, incident management та postmortems.
    \item Інтеграцію з Helm, ArgoCD/GitOps та telemetry stacks.
\end{itemize}

\newpage
\section{Основні принципи SRE}

\subsection{Embracing Risk та Error Budgets}

Надійність — це не 100\% uptime, а баланс між інноваціями та стабільністю.

\begin{itemize}
    \item \textbf{Error Budget}: Дозволений downtime = (1 - SLO). Наприклад, для SLO 99.9\% (monthly) — ~43 хв downtime.
    \item Якщо error budget витрачено — фокус на reliability, а не на нових features.
\end{itemize}

Для ERP/1:
- \texttt{erp-services}: SLO 99.95\%
- \texttt{erp-telemetry}: SLO 99.99\% (критично для observability)

\subsection{SLI — Service Level Indicators}

Приклади SLI для ERP/1 (використовувати Prometheus + Grafana):

\begin{lstlisting}
# Availability
sum(rate(http_requests_total{status!~"5.."}[5m])) / sum(rate(http_requests_total[5m]))

# Latency
histogram_quantile(0.95, sum(rate(http_request_duration_seconds_bucket[5m])) by (le))
\end{lstlisting}

\subsection{Observability Stack (erp-telemetry)}

Використовувати namespaces з \texttt{erpuno/cd/shared} та \texttt{lib/erp-telemetry}:
\begin{itemize}
    \item Prometheus + Alertmanager
    \item Loki / ELK для logs
    \item Jaeger / Tempo для traces
    \item Grafana dashboards per namespace
\end{itemize}

\newpage
\section{Автоматизація та GitOps}

\subsection{Helm + GitOps}

Всі розгортання через Helm charts з \texttt{erpuno/cd/helm}:
\begin{itemize}
    \item \texttt{values.yaml} — конфігурація per environment
    \item \texttt{templates/} — declarative manifests
    \item Shared resources: NetworkPolicies, RBAC, StorageClasses
\end{itemize}

Рекомендації:
\begin{enumerate}
    \item ArgoCD або Flux для GitOps reconciliation
    \item ImagePolicy / Kyverno для security
    \item Helmfile або Helm + Kustomize overlays
\end{enumerate}

\subsection{Continuous Development Pipeline}

Інтеграція SRE в CD:
- Automated testing (unit, integration, chaos)
- Canary / Blue-Green deployments via Helm
- Automated rollbacks on SLO violation

\newpage
\section{Incident Management та On-Call}

\subsection{On-Call Rotations}

\begin{itemize}
    \item PagerDuty / Alertmanager + OpsGenie
    \item Primary + Secondary on-call
    \item Postmortem template (mandatory for all Sev1/Sev2)
\end{itemize}

\subsection{Playbooks}

Створити runbooks в repo для типових інцидентів:
- Pod crashloop
- High latency in erp-ai
- NetworkPolicy violations
- etcd / control-plane issues

\newpage
\section{Chaos Engineering та Resilience}

\begin{itemize}
    \item Chaos Mesh або LitmusChaos в \texttt{erp-infra}
    \item Regular failure injection: node drain, network partition, pod kill
    \item Test recovery time objectives (RTO/RPO)
\end{itemize}

\section{Метрики успіху та TCO}

\begin{itemize}
    \item MTTR < 15 хв
    \item TOIL < 50\% engineer time
    \item Error budget burn rate monitoring
\end{itemize}

\newpage
\section{Висновок}

SRE Handbook ERP/1 є живим документом. Оновлюйте його разом з еволюцією \texttt{erpuno/cd} репозиторію. Головна мета — зробити Continuous Delivery максимально надійною практикою для критичних державних систем.

\begin{flushright}
2026 \copyright\ Zen Crypted
\end{flushright}

\end{document}
